\documentclass[11pt, a4paper]{article}

\usepackage{fontspec}
\setmainfont{Latin Modern Roman}
\setmonofont{DejaVu Sans Mono}[Scale=0.85]
\usepackage{polyglossia}
\setdefaultlanguage{french}
\setotherlanguage{english}
\usepackage[margin=2.5cm]{geometry}
\usepackage{amsmath, amssymb, amsthm}
\usepackage{hyperref}
\usepackage[dvipsnames]{xcolor}
\usepackage{enumitem}
\usepackage{fancyvrb}
\usepackage{caption}
\captionsetup[figure]{name=Listing}
\newcommand{\lean}[1]{\ifmmode\text{\ttfamily\detokenize{#1}}\else{\ttfamily\detokenize{#1}}\fi}
\newcommand{\brk}{\discretionary{}{}{}}
\emergencystretch=5em

\definecolor{indigo}{rgb}{0.29, 0.0, 0.51}
\definecolor{teal}{rgb}{0.0, 0.5, 0.5}

\hypersetup{
    colorlinks=true,
    linkcolor=teal,
    citecolor=indigo,
    urlcolor=blue,
    pdftitle={Une formalisation idiomatique en Mathlib de l'involution de Fricke},
    pdfauthor={Xavier Callens}
}

\newtheorem{theorem}{Théorème}[section]
\newtheorem{lemma}[theorem]{Lemme}
\newtheorem{definition}[theorem]{Définition}
\theoremstyle{remark}
\newtheorem{remark}[theorem]{Remarque}
\renewcommand{\proofname}{Démonstration}

\title{\textbf{Une formalisation idiomatique en Mathlib de \\ l'involution de Fricke et de son opérateur \\ sur les formes modulaires}\\[0.8em]
  \large\textcolor{BrickRed}{PRÉPUBLICATION --- non relue par les pairs}}
\author{
  Xavier Callens\thanks{Auteur correspondant : \texttt{callensxavier@gmail.com}} \\
  \small SocrateAI Research Initiative
}
\date{\today\\[0.6em]
  \small Archivé sur Zenodo : \href{https://doi.org/10.5281/zenodo.22542571}{\texttt{10.5281/zenodo.22542571}} (toutes versions)}

\begin{document}
\maketitle

\begin{center}
\fbox{\begin{minipage}{0.92\textwidth}
\small \textbf{Statut.} Ceci est une \emph{prépublication}. Elle n'a pas fait l'objet d'une
\textbf{relecture par les pairs}. Le développement Lean qu'elle décrit compile et son empreinte
axiomatique est auditée --- ces affirmations sont vérifiées par machine et reproductibles de façon
indépendante --- mais l'\emph{exposition} n'a reçu aucune relecture externe. Cette version intègre
les révisions de deux relectures contradictoires internes : l'une contre une revendication
antérieure d'antériorité mathématique (jugée fausse et rétractée, \S\ref{sec:related}), l'autre
contre une tentative de théorème de pont vers la dualité T (jugée non démontrée et rapportée comme
un résultat négatif, \S\ref{sec:tduality}) ; ni l'une ni l'autre ne constitue une relecture par les
pairs au sens d'une revue. Lire le périmètre avant de citer le résultat.
\end{minipage}}
\end{center}

\begin{abstract}
Nous formalisons, en Lean~4 au-dessus de Mathlib, l'involution de Fricke
$W_N = \left(\begin{smallmatrix} 0 & -1 \\ N & 0 \end{smallmatrix}\right)$ sur $\Gamma_0(N)$
ainsi que l'opérateur qu'elle induit sur les formes modulaires, sous la forme d'une petite API
réutilisable plutôt que d'une infrastructure interne à une démonstration. \textbf{Nous ne
revendiquons aucune antériorité} : du matériel relatif à Fricke et à Atkin--Lehner existe déjà en
Lean~4 dans l'artefact du dernier théorème de Fermat~\cite{FLT2026} (\S\ref{sec:related}), qui
couvre strictement davantage. Ce que nous proposons est un conditionnement destiné à une intégration
en amont : Mathlib elle-même ne contient aucun opérateur de Fricke ni d'Atkin--Lehner. Mathlib
fournit déjà le demi-plan de Poincaré, l'action de Möbius de $GL_2(\mathbb{R})$, l'action
« slash » en poids $k$ portant le facteur déterminant, les sous-groupes de congruence
$\Gamma_0(N)$, ainsi que les types structurés \lean{SlashInvariantForm} et \lean{ModularForm} ;
elle ne contient pas l'opérateur. Nous fournissons cette pièce manquante en trois couches :
(i) au niveau matriciel --- $W_N \in GL_2^+(\mathbb{R})$ avec $\det W_N = N$, l'identité
d'involution $W_N^2 = -N \cdot I$, la conjugaison explicite
$W_N \gamma W_N^{-1} = \left(\begin{smallmatrix} d & -c/N \\ -Nb & a \end{smallmatrix}\right)$,
et le théorème selon lequel $W_N$ normalise $\Gamma_0(N)$ ; (ii) au niveau « slash » ---
$f \mapsto f\vert_k W_N$ préserve l'invariance slash sous $\Gamma_0(N)$ ; et (iii) au niveau
structuré --- l'opérateur est bien défini sur \lean{ModularForm}, son holomorphie et son caractère
borné aux pointes étant transportés le long de $W_N$.
Le développement compte 27 déclarations en 291 lignes ; les 21 théorèmes compilent sans aucun
\lean{sorry} et dépendent exactement des trois axiomes standard de Lean, \lean{propext},
\lean{Classical.choice} et \lean{Quot.sound}, épinglés par 16 contrôles
\texttt{\#guard\_msgs} bloquant la compilation, de sorte qu'une dérive de l'empreinte axiomatique
provoque une erreur de compilation. Nous ne revendiquons aucune mathématique nouvelle : les
résultats sont classiques, dus à Fricke et à Atkin--Lehner~\cite{AtkinLehner1970}. L'apport est la
mécanisation, et son usage visé est de servir d'infrastructure à la formalisation des identités de
quotients êta, où les valeurs propres de Fricke constituent l'invariant organisateur.
\end{abstract}

\noindent\textbf{MSC 2020 :} 11F11, 11F06, 68V20.\quad
\textbf{Mots-clés :} involution de Fricke, théorie d'Atkin--Lehner, sous-groupes de congruence,
Lean~4, Mathlib, vérification formelle.

\section{Introduction}

L'involution de Fricke s'énonce de façon élémentaire et intervient partout dans la théorie des
formes modulaires sur $\Gamma_0(N)$. Pour $N \ge 1$, posons
\[
  W_N \;=\; \begin{pmatrix} 0 & -1 \\ N & 0 \end{pmatrix}, \qquad \det W_N = N > 0 .
\]
Agissant sur le demi-plan de Poincaré par $\tau \mapsto -1/(N\tau)$, elle échange les pointes $0$
et $\infty$ de $\Gamma_0(N)$, normalise $\Gamma_0(N)$, et induit une involution de
$M_k(\Gamma_0(N))$ dont les sous-espaces propres $\pm 1$ organisent une large part de
l'arithmétique de l'espace --- de la manière la plus visible pour les quotients êta, où la
multiplicativité se lit sur les valeurs propres de Fricke et
d'Atkin--Lehner~\cite{Martin1996, Ono2004}.

Cette note est écrite dans le mode que la pratique récente rend standard : un artefact formel livré
conjointement à l'exposition destinée aux lecteurs humains. L'autoformalisation à grande échelle a
franchi un seuil --- une démonstration complète et vérifiée par machine du dernier théorème de
Fermat a été produite en onze jours par des agents d'intelligence artificielle coordonnés par un
graphe d'énoncés partagé~\cite{FLT2026, Prove2Me2026} --- et les questions méthodologiques
afférentes relèvent désormais des mathématiques ordinaires~\cite{Tao2026AI, Buzzard2021}. Nous
travaillons en Lean~4~\cite{deMoura2021} au-dessus de Mathlib~\cite{Mathlib2020}, et nous adoptons
deux pratiques précises issues de l'artefact FLT : l'audit axiomatique en tant que \emph{cible de
compilation}, et un graphe explicite de dépendances entre énoncés (\S\ref{sec:verification}).

\subsection*{Ce que Mathlib fournit déjà}

Nous sommes explicites sur l'état de l'art, car il détermine ce qui est nouveau ici et ce qui ne
l'est pas. Au commit \texttt{905b9581} (28 juillet 2026), Mathlib~\cite{Mathlib2020} contient, sous
des formes strictement plus générales que ce qu'il nous faudrait réénoncer :

\begin{itemize}
  \item \lean{UpperHalfPlane} et l'action de Möbius de \lean{GL (Fin 2) ℝ}, ainsi que
        \lean{denom_ne_zero_of_im} et \lean{normSq_denom_pos} ;
  \item \lean{im_smul_eq_div_normSq}, qui est exactement l'identité classique
        $\operatorname{Im}(M \cdot \tau) = |\det M| \operatorname{Im}(\tau) / |c\tau + d|^2$,
        énoncée pour $M \in GL_2(\mathbb{R})$ arbitraire via $|\det M|$ et non seulement pour
        $\det M > 0$ ;
  \item \lean{Matrix.GLPos}, le sous-groupe des matrices de déterminant positif ;
  \item \lean{CongruenceSubgroup.Gamma0 N} ;
  \item l'action « slash » en poids $k$, \lean{SlashAction ℤ (GL (Fin 2) ℝ) (ℍ → ℂ)}, dont la
        définition \emph{porte déjà} le facteur déterminant :
        $f\vert_k g = \sigma(g)\, f(g\cdot\tau)\, |\det g|^{k-1} \operatorname{denom}(g,\tau)^{-k}$ ;
  \item le type structuré \lean{ModularForm Γ k} pour
        \lean{Γ : Subgroup (GL (Fin 2) ℝ)}, avec holomorphie et caractère borné aux pointes.
\end{itemize}

Rien de tout cela n'est redémontré ici ; nous l'importons. Un lecteur cherchant une formalisation
de l'action de $GL_2^+(\mathbb{R})$ sur $\mathbb{H}$ doit utiliser celle de Mathlib, non la nôtre.

\subsection*{La lacune, et notre apport}

Une recherche sur l'ensemble du dépôt Mathlib à ce commit pour \texttt{Fricke} ou
\texttt{AtkinLehner} ne renvoie \emph{aucune} occurrence : la machinerie permettant d'héberger ces
opérateurs existe dans Mathlib, mais les opérateurs n'y sont pas. C'est la lacune que nous comblons
--- avec la réserve importante, développée au \S\ref{sec:related}, que du matériel relatif à Fricke
\emph{existe bel et bien} en Lean~4 hors de Mathlib. Nos apports sont donc énoncés comme une API,
et non comme des premières :

\begin{enumerate}
  \item[(C1)] $W_N$ comme élément de \lean{GL (Fin 2) ℝ} avec $\det W_N = N$, et
        $W_N \in GL_2^+(\mathbb{R})$ (théorème~\ref{thm:glpos}) ;
  \item[(C2)] $W_N^2 = -N \cdot I$ (théorème~\ref{thm:involution}), ce qui, les matrices scalaires
        agissant trivialement sur $\mathbb{H}$, constitue la propriété d'involution ;
  \item[(C3)] la formule de conjugaison explicite (théorème~\ref{thm:conj}) ;
  \item[(C4)] $W_N$ conjugue $\Gamma_0(N)$ dans lui-même à l'intérieur de $GL_2(\mathbb{R})$
        (théorème~\ref{thm:normalizes}) ; joint au fait que $W_N^2 = -N\cdot I$ est central, cela
        donne la condition de normalisateur bilatère, mais c'est l'inclusion que nous formalisons ;
  \item[(C5)] le slash par $W_N$ préserve l'invariance slash sous $\Gamma_0(N)$, d'où l'opérateur
        sur \lean{SlashInvariantForm} (théorème~\ref{thm:slashinv}) ;
  \item[(C6)] l'holomorphie et le caractère borné aux pointes se transportent le long de $W_N$,
        d'où l'opérateur sur le type structuré \lean{ModularForm} (théorème~\ref{thm:modular}) ;
  \item[(C7)] le composé $f \mapsto (f\vert_k W_N)\vert_k W_N$ est la multiplication par le
        scalaire explicite $(-1)^k N^{k-2}$ (théorème~\ref{thm:composite}), ce qui identifie la
        normalisation faisant de $\vert_k W_N$ une involution.
\end{enumerate}

(C4) est le point charnière : c'est ce qui rend $f \mapsto f\vert_k W_N$ bien défini. (C5) et (C6)
le remontent jusqu'aux types structurés de Mathlib, afin que l'opérateur de Fricke puisse se
composer avec le reste de la bibliothèque plutôt que de vivre à côté d'elle.

\section{Rappels mathématiques}

Soit $\gamma = \left(\begin{smallmatrix} a & b \\ c & d\end{smallmatrix}\right) \in SL_2(\mathbb{Z})$
avec $\gamma \in \Gamma_0(N)$, c'est-à-dire $N \mid c$ ; écrivons $c = Nc'$. Un calcul direct donne
\begin{equation}\label{eq:conj}
  W_N \gamma W_N^{-1}
  \;=\; \begin{pmatrix} d & -c' \\ -Nb & a \end{pmatrix} ,
\end{equation}
dont les coefficients sont entiers précisément parce que $N \mid c$, dont le déterminant vaut
$da - (-c')(-Nb) = ad - bc = 1$, et dont le coefficient inférieur gauche $-Nb$ est divisible par
$N$. Le conjugué appartient donc de nouveau à $\Gamma_0(N)$, et $W_N$ normalise $\Gamma_0(N)$.

Plutôt que de manipuler $W_N^{-1}$ dans la formalisation, nous démontrons l'identité
\emph{d'entrelacement} équivalente
\begin{equation}\label{eq:intertwine}
  W_N \cdot \gamma \;=\; \delta \cdot W_N, \qquad
  \delta := \begin{pmatrix} d & -c' \\ -Nb & a \end{pmatrix},
\end{equation}
qui évite entièrement les inverses et dont \eqref{eq:conj} découle par une multiplication à droite.
Les deux membres de \eqref{eq:intertwine} valent
$\left(\begin{smallmatrix} -c & -d \\ Na & Nb \end{smallmatrix}\right)$, l'accord ayant lieu
exactement sous la relation $c = Nc'$.

\section{Le développement Lean}

Le développement comporte quatre modules. Nous reproduisons les déclarations porteuses verbatim
depuis les sources compilées.

\begin{figure}[h]
\begin{Verbatim}[frame=single,fontsize=\small,samepage=true]
def frickeMatrix (N : ℕ) : Matrix (Fin 2) (Fin 2) ℝ := !![0, -1; (N : ℝ), 0]

@[simp] theorem frickeMatrix_det (N : ℕ) : (frickeMatrix N).det = (N : ℝ) := by
  simp [frickeMatrix, Matrix.det_fin_two_of]

noncomputable def frickeW {N : ℕ} (hN : 0 < N) : GL (Fin 2) ℝ :=
  Matrix.GeneralLinearGroup.mkOfDetNeZero (frickeMatrix N) (frickeMatrix_det_ne_zero hN)
\end{Verbatim}
\caption{La matrice de Fricke et son déterminant.}
\end{figure}

\begin{theorem}[$W_N$ est de déterminant positif]\label{thm:glpos}
\lean{frickeW_mem_GLPos : frickeW hN ∈ GLPos (Fin 2) ℝ}.
\end{theorem}

\begin{theorem}[Involution]\label{thm:involution}
$W_N^2 = -N \cdot I$ :
\lean{frickeW_sq_coe : ((frickeW hN * frickeW hN : GL (Fin 2) ℝ) : Matrix _ _ ℝ)}
\lean{= (-(N : ℝ)) • 1}.
\end{theorem}

Les matrices scalaires agissant trivialement sur $\mathbb{H}$, le théorème~\ref{thm:involution}
exprime précisément que $W_N$ induit une involution du demi-plan de Poincaré.

L'ingrédient arithmétique est la divisibilité extraite de l'appartenance à $\Gamma_0(N)$, que
Mathlib énonce modulo $N$ :

\begin{figure}[h]
\begin{Verbatim}[frame=single,fontsize=\small,samepage=true]
theorem gamma0_dvd_lower_left {N : ℕ} {γ : SL(2, ℤ)} (hγ : γ ∈ Gamma0 N) :
    (N : ℤ) ∣ γ.1 1 0 := by
  rw [Gamma0_mem] at hγ
  exact (ZMod.intCast_zmod_eq_zero_iff_dvd _ _).mp hγ
\end{Verbatim}
\caption{Divisibilité du coefficient inférieur gauche.}
\end{figure}

Le conjugué $\delta$ de \eqref{eq:conj} se définit alors comme un véritable élément de
$SL_2(\mathbb{Z})$, sa condition de déterminant étant acquittée par \lean{linear_combination}
contre $\det \gamma = 1$ :

\begin{figure}[h]
\begin{Verbatim}[frame=single,fontsize=\footnotesize,samepage=true]
def frickeConj {N : ℕ} (γ : SL(2, ℤ)) (c' : ℤ) (hc : γ.1 1 0 = (N : ℤ) * c') : SL(2, ℤ) :=
  ⟨!![γ.1 1 1, -c'; -(N : ℤ) * γ.1 0 1, γ.1 0 0], by
    have hdet := γ.2
    rw [Matrix.det_fin_two] at hdet; rw [hc] at hdet
    rw [Matrix.det_fin_two_of]; linear_combination hdet⟩
\end{Verbatim}
\caption{Le conjugué de Fricke, avec sa preuve de déterminant.}
\end{figure}

\begin{theorem}[Conjugaison]\label{thm:conj}
\lean{frickeW_conj_eq} : pour $\gamma \in SL_2(\mathbb{Z})$ avec $c = Nc'$,
\[
  W_N \cdot \mathrm{mapGL}_{\mathbb{R}}(\gamma) \cdot W_N^{-1}
  \;=\; \mathrm{mapGL}_{\mathbb{R}}(\mathrm{frickeConj}\,\gamma\,c') .
\]
La démonstration réécrit par l'identité d'entrelacement \eqref{eq:intertwine} et simplifie
$W_N W_N^{-1}$.
\end{theorem}

\begin{theorem}[Résultat principal : normalisation]\label{thm:normalizes}
\lean{frickeW_normalizes_Gamma0} :
\[
  \forall g \in \mathrm{Gamma0}(N).\mathrm{map}(\mathrm{mapGL}_{\mathbb{R}}), \quad
  W_N \, g \, W_N^{-1} \in \mathrm{Gamma0}(N).\mathrm{map}(\mathrm{mapGL}_{\mathbb{R}}) .
\]
\end{theorem}

\begin{proof}[Démonstration (telle que formalisée)]
Décomposer $g = \mathrm{mapGL}_{\mathbb{R}}(\gamma)$ avec $\gamma \in \Gamma_0(N)$ ; obtenir $c'$
tel que $c = Nc'$ par \lean{gamma0_dvd_lower_left} ; réécrire par le théorème~\ref{thm:conj} ;
conclure par \lean{Subgroup.mem_map_of_mem} appliqué à \lean{frickeConj_mem_Gamma0}, dont le contenu
est que le coefficient inférieur gauche $-Nb$ s'annule modulo $N$.
\end{proof}

\section{L'opérateur sur les formes modulaires}

La normalisation est un énoncé sur les groupes ; être un \emph{opérateur sur les formes modulaires}
exige de transporter le long de $W_N$ les deux conditions analytiques de \lean{ModularForm}. Nous
procédons en deux étapes, en travaillant avec l'image
$\Gamma_0(N)_{\mathbb{R}} := \Gamma_0(N).\mathrm{map}(\mathrm{mapGL}_{\mathbb{R}})$ dans
$GL_2(\mathbb{R})$, qui est le sous-groupe sur lequel les types structurés de Mathlib sont énoncés.

\begin{theorem}[Invariance slash]\label{thm:slashinv}
\lean{slash_frickeW_invariant} : pour $f$ une \lean{SlashInvariantForm} de poids $k$ pour
$\Gamma_0(N)_{\mathbb{R}}$ et tout $g$ dans ce sous-groupe,
$(f\vert_k W_N)\vert_k g = f\vert_k W_N$.
\end{theorem}

\begin{proof}[Démonstration (telle que formalisée)]
Écrire $g = \mathrm{mapGL}_{\mathbb{R}}(\gamma)$ avec $\gamma \in \Gamma_0(N)$ et extraire
$c = Nc'$. Alors, en utilisant le cocycle du slash dans les deux sens et l'identité
d'entrelacement~\eqref{eq:intertwine},
\[
  (f\vert_k W_N)\vert_k \gamma \;=\; f\vert_k(W_N\gamma) \;=\; f\vert_k(\delta W_N)
  \;=\; (f\vert_k \delta)\vert_k W_N \;=\; f\vert_k W_N ,
\]
la dernière étape par invariance slash de $f$ sous $\delta \in \Gamma_0(N)$
(théorème~\ref{thm:conj}). Cela se conditionne en \lean{frickeSlashOperator}.
\end{proof}

\begin{theorem}[Opérateur sur \lean{ModularForm}]\label{thm:modular}
\lean{frickeModularOperator} : $f \mapsto f\vert_k W_N$ est une application
$\lean{ModularForm}\ \Gamma_0(N)_{\mathbb{R}}\ k \to \lean{ModularForm}\ \Gamma_0(N)_{\mathbb{R}}\ k$.
\end{theorem}

\begin{proof}[Démonstration (telle que formalisée)]
L'holomorphie est immédiate à partir de \lean{MDifferentiable.slash} de Mathlib : le slash par tout
élément fixé de $GL_2(\mathbb{R})$ préserve \lean{MDiff}. Pour le caractère borné,
\lean{OnePoint.IsBoundedAt.smul_iff} de Mathlib réduit \emph{$f\vert_k W_N$ est borné en $c$} à
\emph{$f$ est borné en $W_N \cdot c$}, de sorte qu'il suffit que $W_N$ envoie les pointes de
$\Gamma_0(N)$ sur des pointes de $\Gamma_0(N)$. C'est \lean{isCusp_frickeW_smul}, obtenu à partir
de \lean{IsCusp.smul} de Mathlib, de la forme « sous-groupe » du théorème~\ref{thm:normalizes}
(\lean{frickeW_conjAct_le}) et de la monotonie de \lean{IsCusp} en son argument de sous-groupe.
\end{proof}

Concrètement, $W_N$ échange les pointes $0$ et $\infty$ de $\Gamma_0(N)$ ; ce qui est démontré ici
est l'énoncé plus faible, suffisant pour le transport, que l'ensemble des pointes est préservé.

\begin{theorem}[Le composé est un scalaire]\label{thm:composite}
\lean{frickeW_sq_slash} : pour toute $f : \mathbb{H} \to \mathbb{C}$ et tout poids $k$,
\[
  (f\vert_k W_N)\vert_k W_N \;=\; (N^2)^{k-1}(-N)^{-k}\, f \;=\; (-1)^k N^{k-2}\, f .
\]
\end{theorem}

\begin{proof}[Démonstration (telle que formalisée)]
Par la loi de cocycle, le membre de gauche vaut $f\vert_k(W_N^2)$, et $W_N^2 = -N\cdot I$ est
scalaire (théorème~\ref{thm:involution}). Pour cette matrice, les trois ingrédients du slash de
Mathlib se calculent directement : $\det(W_N^2) = N^2$, $\operatorname{denom}(W_N^2,\tau) = -N$, et
$W_N^2$ agit trivialement sur $\mathbb{H}$. Le facteur de conjugaison $\sigma$ ne nécessite aucune
disjonction de cas sur le déterminant : $\sigma(gg') = \sigma g \circ \sigma g'$ et
$\sigma g \circ \sigma g = \mathrm{id}$ donnent $\sigma(W_N^2) = \mathrm{id}$ immédiatement.
\end{proof}

Pour $k$ pair, il s'ensuit que $f \mapsto N^{1-k/2}(f\vert_k W_N)$ est de carré l'identité --- c'est
l'involution de Fricke proprement dite. Pour $k$ impair la question est vide, puisque
$-I \in \Gamma_0(N)$ force $f \equiv 0$ (\lean{ModularForm.eq_zero_of_neg_one_mem} de Mathlib).

\section{État de la vérification}\label{sec:verification}

Toutes les déclarations du module compilent. Il n'y a aucun \lean{sorry} et aucun axiome propre au
projet. L'empreinte axiomatique de chaque théorème, obtenue par \texttt{\#print axioms}, est
exactement
\[
  [\,\texttt{propext},\ \texttt{Classical.choice},\ \texttt{Quot.sound}\,] ,
\]
les trois axiomes standard de la logique de Lean. À la suite de l'artefact du dernier théorème de
Fermat~\cite{FLT2026}, l'empreinte n'est pas seulement rapportée mais \emph{imposée} : un module
\lean{FinalCheck} place une garde \texttt{\#guard\_msgs in \#print axioms} sur chacun des seize
théorèmes principaux à l'intérieur de la cible de compilation par défaut, de sorte que tout
élargissement de l'ensemble d'axiomes est une erreur de compilation ; un contrôle négatif (une
garde affirmant une empreinte délibérément fausse) est conservé hors de la compilation et il a été
vérifié qu'il échoue. Le développement est suivi comme un graphe de dépendances entre énoncés
(\lean{dag/theorems.jsonl}, avec une carte lisible dans \lean{dag/PROOF-PATH.md}) ; la machine
vérifie que tout n\oe{}ud marqué démontré nomme une déclaration existante et ne dépend que de
n\oe{}uds démontrés. Nous ne revendiquons rien au-delà de ce que le noyau vérifie : dans la
terminologie adoptée dans l'ensemble du programme SocrateAI, (C1)--(C7) sont des résultats
\emph{formels} ; l'exposition qui les entoure ne l'est pas.

\paragraph{Artefact.}
Archivé sur Zenodo, \href{https://doi.org/10.5281/zenodo.22542571}{\texttt{doi:10.5281/zenodo.22542571}},
qui constitue la version de référence citable de cette note et de ses sources.
Dépôt \url{https://github.com/xaviercallens/SocrateAI-Lean-Lib} (étiquette \texttt{fricke-v1}) ;
modules
\texttt{Lean/\brk SocrateAI/\brk ModularForms/\brk \{FrickeInvolution,\brk\ FrickeSlash,\brk\ FrickeModular,\brk\ FrickeComposite\}.lean},
les gardes axiomatiques se trouvant dans \texttt{Lean/\brk SocrateAI/\brk FinalCheck.lean} ;
un miroir accompagné d'une fiche descriptive est disponible à
\url{https://huggingface.co/datasets/callensxavier/socrateai-fricke-involution}.
Chaîne d'outils \texttt{leanprover/lean4:v4.32.2} ;
Mathlib figée au commit \texttt{905b9581} (28 juillet 2026) ;
\texttt{lake build SocrateAI} s'achève en 3053 tâches (module de gardes axiomatiques inclus).
Les quatre modules totalisent 27 déclarations en 291 lignes et n'importent que
\texttt{Mathlib.\allowbreak NumberTheory.\allowbreak ModularForms.\allowbreak SlashActions},
\texttt{Mathlib.\allowbreak NumberTheory.\allowbreak ModularForms.\allowbreak CongruenceSubgroups},
\texttt{Mathlib.\allowbreak NumberTheory.\allowbreak ModularForms.\allowbreak SlashInvariantForms} et
\texttt{Mathlib.\allowbreak NumberTheory.\allowbreak ModularForms.\allowbreak Basic}.
La configuration de compilation est distribuée et portable : \texttt{lakefile.lean} requiert Mathlib
depuis git à la révision indiquée ci-dessus, \texttt{lean-toolchain} lui correspond, et aucun fichier
suivi ne contient de chemin propre à une machine. \texttt{git clone}, \texttt{lake exe cache get},
\texttt{lake build SocrateAI} reproduit l'artefact ; voir \texttt{BUILDING.md}. Les versions
antérieures de cette note indiquaient que la configuration n'était pas distribuée. C'était exact, et
en deçà de la réalité : le \texttt{lakefile.lean} publié ne déclarait aucune dépendance à Mathlib et
son \texttt{lean-toolchain} nommait une autre version du compilateur, de sorte qu'un clone propre ne
pouvait pas compiler. Les deux points sont corrigés.

\section{Travaux apparentés}\label{sec:related}

\paragraph{Formalisation antérieure en Lean~4 de matériel relatif à Fricke.}
L'artefact du dernier théorème de Fermat~\cite{FLT2026} est un développement Lean~4 construit
sur Mathlib qui contient déjà une théorie substantielle de Fricke et d'Atkin--Lehner. Une recherche
de code à l'échelle du dépôt renvoie de l'ordre de $2{,}7\times10^3$ fichiers Lean mentionnant
\texttt{Fricke} et $3{,}6\times10^3$ mentionnant \texttt{AtkinLehner}, dont
\texttt{Thm\_ModularForm\_heckeU\_add\brk\_slash\brk\_fricke\brk\_eq\brk\_zero},
\texttt{Thm\_CuspForm\_hasNebentypus\brk\_inv\brk\_and\brk\_qCoeff\brk\_hecke\brk\_eigen\brk\_of\brk\_fricke} et
\texttt{Def\_ModularCurve\_AtkinLehner}. Leurs énoncés utilisent la même matrice
$\left(\begin{smallmatrix} 0 & -1 \\ N & 0\end{smallmatrix}\right)$ et la même action slash de
Mathlib que nous employons ici, et ils vont considérablement plus loin --- jusqu'aux opérateurs de
Hecke, aux formes propres à nebentypus, et aux opérateurs d'Atkin--Lehner sur les courbes
modulaires.

De la façon la plus directe, \texttt{Def\_CuspForm\_AtkinLehnerOperator.lean} définit
\texttt{atkinLehnerLin}, une application
$\texttt{ModularForm}\ (\Gamma_0(M))\ k \to_{\ell[\mathbb{C}]} \texttt{ModularForm}\ (\Gamma_0(M))\ k$,
acquittant l'invariance slash, l'holomorphie et le caractère borné aux pointes comme le fait notre
\texttt{frickeModularOperator} --- soit nos (C5) et (C6), pour toute la famille d'Atkin--Lehner et
sous forme d'application $\mathbb{C}$-linéaire. Leur
\texttt{Def\_ModularForm\_AtkinLehnerDatum.lean} construit son élément de groupe avec le même
idiome \texttt{mkOfDetNeZero} que nous et démontre l'analogue de nos (C1), (C2) et (C7).

En particulier, leur \texttt{CohCarrier.frickeMat} est, à la dénomination près, notre
\texttt{frickeConj} : la même matrice conjuguée
$\left(\begin{smallmatrix} d & -c/N \\ -Nb & a\end{smallmatrix}\right)$, la même hypothèse de
divisibilité, et une preuve de déterminant qui s'achève elle aussi par
\texttt{linear\_combination}. Notre développement a été obtenu indépendamment, ce qui n'a aucune
incidence sur l'antériorité ; nous consignons la coïncidence parce qu'un lecteur comparant les deux
doit savoir qu'il s'agit de la même construction.

\paragraph{Ce qui est néanmoins proposé ici.}
Le matériel FLT est interne à une démonstration : il vit dans un espace de noms \texttt{CohCarrier}
ad hoc, et $W_N$ y est généralement passé à chaque théorème comme un paramètre contraint par une
équation matricielle. Nous définissons au contraire $W_N$ une fois pour toutes comme un élément de
\texttt{GL (Fin 2) ℝ} et démontrons la petite théorie environnante --- appartenance à $GL_2^+$,
$W_N^2 = -N\cdot I$, la formule de conjugaison, et les opérateurs induits sur
\texttt{SlashInvariantForm} et \texttt{ModularForm}. La destination visée est Mathlib, qui ne
contient rien de tout cela. Énoncé précisément, ce qui subsiste après divulgation est : la matrice
littérale $W_N$ sur $\mathbb{R}$ (le coefficient inférieur droit de leur donnée est $q$, si bien
qu'elle ne se spécialise pas exactement en $W_N$), le scalaire composé explicite de (C7), et un
développement à quatre importations nommées utilisant les tactiques standard plutôt qu'un harnais
de preuve dédié. C'est une différence de méthode, et savoir si elle vaut quelque chose est une
question pour la relecture Mathlib, non une chose qu'il nous appartient d'affirmer.

\section{Limites}

Les opérateurs de Fricke et d'Atkin--Lehner sont classiques ; le traitement de référence des
involutions d'Atkin--Lehner $W_Q$ pour $Q \| N$ est~\cite{AtkinLehner1970}, et des exposés de
manuel figurent dans~\cite{DiamondShurman2005, Miyake2006, Ono2004}. Notre apport n'est pas une
nouveauté mathématique mais le transfert de ce matériel dans un cadre vérifié par machine où l'on
peut s'appuyer dessus.

Nous énonçons les limites sans détour.

\begin{remark}[Ce qui n'est \emph{pas} fait ici]
\leavevmode
\begin{enumerate}[label=(\alph*)]
  \item Nous ne formalisons que l'involution de Fricke $W_N$. La famille complète d'Atkin--Lehner
        $W_Q$ pour $Q \| N$, et la décomposition simultanée en sous-espaces propres qui en résulte,
        ne sont pas traitées.
  \item La décomposition de $M_k(\Gamma_0(N))$ en sous-espaces propres $\pm 1$ n'est pas elle-même
        conditionnée comme une définition Lean. Le contenu mathématique est le
        théorème~\ref{thm:composite}, qui est formalisé ; ce qui reste relève de la tenue de
        comptes --- porter l'hypothèse « $k$ pair » à travers l'arithmétique \lean{zpow} sur la
        constante, et empaqueter les deux sous-espaces propres comme sous-modules.
  \item Nous ne faisons aucune revendication d'antériorité. Notre vérification de lacune a été
        menée contre Mathlib, où elle est valide ; une recherche ultérieure dans les développements
        Lean~4 cités ici a mis au jour l'état de l'art décrit au \S\ref{sec:related}. Nous n'avons
        exploré ni le Zulip Lean ni les autres assistants de preuve.
\end{enumerate}
\end{remark}

\section{Application visée}

La motivation est en aval. Les quotients êta
$f(\tau) = \prod_{\delta \mid N} \eta(\delta \tau)^{r_\delta}$ sont modulaires sur $\Gamma_0(N)$
sous les conditions de Ligozat, et leur comportement sous les opérateurs de Fricke et
d'Atkin--Lehner est ce qui rend la multiplicativité traitable~\cite{Martin1996}. Une formalisation
des identités de quotients êta requiert donc $W_N$ comme prérequis, et c'est ce que ce module
fournit. Une première tentative de préciser ce lien est rapportée au \S\ref{sec:tduality} ci-dessous, et
elle a échoué de façon instructive plutôt que de réussir. L'opérateur étant désormais défini sur \lean{ModularForm}, les
étapes suivantes immédiates sont la constante de normalisation de la limite (b) et la famille
complète d'Atkin--Lehner de la limite (a). Du côté analytique, les bornes de coefficients de type
Rademacher qui accompagnent le programme des quotients êta seront d'abord vérifiées au niveau
« validé par programme » à l'aide d'un outillage léger de preuve d'inégalités dans l'esprit de
l'assistant \emph{estimates} de Tao~\cite{TaoEstimates2025}, et seulement ensuite promues au rang
d'énoncés vérifiés par le noyau --- la même discipline à deux niveaux (vérificateur exact en
dessous, noyau au-dessus) employée dans tout ce programme.

\section{Un résultat négatif : une tentative de pont vers la dualité T}\label{sec:tduality}

Un travail de suivi, rapporté ici pour mémoire plutôt que laissé comme une simple anticipation non
qualifiée, a tenté de formaliser que $W_N$ n'est pas seulement un outil pour les quotients êta mais
l'incarnation en Lean~4 d'une dualité physique : que l'action de la dualité T sur le module de
structure complexe $\tau$ d'un facteur $T^2$, restreinte à $\Gamma_0(N)$, \emph{coïncide} avec
l'action de $W_N$ sur $\mathbb{H}$, déjà démontrée. Cette section énonce clairement ce qui a été
trouvé, car taire un résultat négatif est précisément le mode de défaillance que ce programme a déjà
corrigé deux fois auparavant~\cite{FLT2026} --- et parce que ce qui survit à la tentative est
authentique et réutilisable.

\subsection*{Ce qui est démontré, sans \lean{sorry}, et réutilisable}

\begin{theorem}[L'action de Möbius de $W_N$, sous forme close]\label{thm:tdualm1}
\lean{frickeW_smul_coe} : pour tout $N>0$ et $\tau\in\mathbb{H}$,
\[
  \bigl(W_N \cdot \tau : \mathbb{C}\bigr) \;=\; \frac{-1}{N\tau}.
\]
\end{theorem}

Ceci referme une véritable lacune que rien dans la bibliothèque ne refermait auparavant : la forme
normale $-1/(N\tau)$ est exactement ce que la moitié à valeurs dans $\mathbb{C}$ des quotients êta
de ce programme calcule déjà (\lean{eta_fricke_factor}, dans les modules de quotients êta), alors
que $W_N$ lui-même n'avait jamais été énoncé qu'au langage de l'action de $GL_2(\mathbb{R})$. La
preuve passe par la disjonction de cas sur le déterminant de Mathlib (le théorème~\ref{thm:glpos}
fournit exactement l'hypothèse de déterminant positif dont le lemme d'action propre de Mathlib a
besoin, si bien qu'aucun signe ni conjugaison cachés n'entrent en jeu) et est certifiée non vide par
un contrôle négatif : le même script de preuve appliqué à quatre membres de droite \emph{faux},
calculés indépendamment ($+1/(N\tau)$, $-1/\tau$, $-N/\tau$, $-1/(N+\tau)$), échoue avec quatre
\lean{unsolved goals} distincts, chacun affichant le même membre de gauche correct.

\subsection*{Ce qui n'est pas démontré, et pourquoi}

L'énoncé du pont lui-même,
\[
  \bigl(W_N \cdot \tau : \mathbb{C}\bigr) = \tfrac{-1}{N\tau}
  \ \wedge\ \forall\,\gamma\in\Gamma_0(N),\ \gamma\cdot\tau
    \;\text{est l'action fractionnaire-linéaire standard de}\ SL_2(\mathbb{Z}),
\]
reste \lean{sorry}, sous une garde d'axiomes inversée qui fait échouer la construction, et une
relecture contradictoire du propre déroulement de ce programme a trouvé deux raisons indépendantes
pour qu'il en reste ainsi.

\textbf{C'est un renommage, pas une réduction.} Un contrôle d'équivalence autonome, compilé et
vérifié, montre que l'énoncé du pont est logiquement équivalent à la conjonction de trois lemmes
déjà démontrés et n'ajoute aucun contenu au-delà d'eux : l'hypothèse $\gamma\in\Gamma_0(N)$ n'est
jamais consommée (le second conjoint vaut pour tout $\gamma\in SL_2(\mathbb{Z})$), et ni le
théorème~\ref{thm:normalizes} (que $W_N$ normalise $\Gamma_0(N)$) ni la propriété d'involution
$W_N^2=-N\cdot I$ n'apparaissent nulle part dans l'énoncé ou son terme de preuve. Rien dans le terme
Lean ne nomme un groupe de dualité, un fond de théorie des cordes, ou un spectre ; la physique n'est
affirmée que dans la prose autour du théorème, pas à l'intérieur.

\textbf{La référence en physique ne le soutient pas.} La référence naturelle pour une action de
dualité T sur un module $\tau$, Giveon--Porrati--Rabinovici~\cite{GPR1994}, n'ancre
$\tau\mapsto-1/\tau$ qu'à $N=1$. Pour $N>1$, $\det W_N = N$, si bien que $W_N$ n'est même pas un
élément du $SL_2(\mathbb{Z})$ que décrit cette référence, et une recherche exhaustive de son texte
ne retourne aucune occurrence d'une involution de Fricke, d'un sous-groupe de congruence, ou d'un
élément de dualité dépendant du niveau, pour un $N>1$ quelconque. Le foyer physique correct pour une
application de type Fricke de niveau $N$ s'avère être un objet tout autre :
Persson--Volpato~\cite{PerssonVolpato2015} identifient $S\mapsto -1/(NS)$ comme une
\emph{dualité S} sur le module axio-dilaton hétérotique dans les modèles orbifold CHL,
explicitement \emph{en dehors} de la symétrie standard $SL_2(\mathbb{Z})$ de la théorie parente ---
un module différent, et une dualité différente, de ceux visés par cette tentative.

\subsection*{La conclusion honnête}

Le théorème~\ref{thm:tdualm1} est réel : c'est une preuve que $W_N$ est le bon objet mathématique
pour une correspondance physique, énoncé purement comme un fait sur $\mathbb{H}$ et $\mathbb{C}$
sans aucun contenu physique revendiqué. Quelle correspondance, et sous quelles hypothèses, reste
ouvert, et la référence de Persson--Volpato est le point de départ de la prochaine tentative ---
à l'intérieur des orbifolds CHL de $K3\times T^2$, sur l'axio-dilaton, non le module de $T^2$. Nous
consignons ceci comme une conclusion plutôt que d'abandonner discrètement la phrase qui le
promettait.

\section{Conclusion}

Nous avons formalisé en Lean~4 l'involution de Fricke $W_N$ sur $\Gamma_0(N)$ et l'opérateur
qu'elle induit sur les formes modulaires, comblant une lacune de \emph{Mathlib} --- mais non, comme
le consigne le \S\ref{sec:related}, de Lean~4 dans son ensemble. Le développement est petit ---
27 déclarations en 291 lignes ---, ne dépend que des trois axiomes standard de Lean, et est bâti
sur l'infrastructure \lean{SlashAction}, \lean{SlashInvariantForm} et \lean{ModularForm} existante
de Mathlib plutôt que de la dupliquer, de sorte qu'il se compose avec le reste de la bibliothèque.
Les mathématiques sont classiques ; leur énoncé vérifié par machine ne l'est pas.

\section*{Note sur cette version}

Le présent texte est la version française d'une note rédigée initialement en anglais et déposée sur
Zenodo sous le DOI \texttt{10.5281/zenodo.22542571}. Les deux versions ont le même contenu
mathématique ; en cas de divergence, la version anglaise fait foi pour les énoncés formels, dont
les noms Lean ne sont pas traduits. Cette note n'a pas fait l'objet d'une relecture par les pairs :
notre propre relecture interne a conclu qu'elle ne constituait pas une soumission de revue, et la
voie recommandée pour ce matériel est une contribution (\emph{pull request}) à Mathlib. Le dépôt
HAL vaut archivage et citation, non validation.

\section*{Note sur la méthode : assistance de l'IA et responsabilité humaine}

Cette note suit la norme de transparence défendue par Tao~\cite{Tao2026AI}, qui reprend les
recommandations de la Déclaration de Leyde sur l'intelligence artificielle et les
mathématiques~\cite{Leiden2026} : divulguer l'usage d'outils, affirmer que le crédit et la
responsabilité reviennent à l'auteur humain, et attribuer les sources là où les outils d'IA ne
peuvent pas le faire de manière fiable eux-mêmes.

\textbf{Ce que l'IA a fait.} La formalisation en Lean~4, la recherche bibliographique et la
rédaction de ce texte ont été produites avec l'assistance de grands modèles de langage (la famille
Claude d'Anthropic), sous direction humaine sur de nombreuses sessions. La méthodologie est
neuro-symbolique au sens précis que décrit Tao~\cite{Tao2026AI} : le raisonnement en langage
naturel propose des définitions, des stratégies de preuve et de la prose, et c'est le noyau Lean ---
non le modèle, ni la confiance de l'auteur dans le modèle --- qui admet ou rejette chaque
affirmation mathématique. Chaque théorème de cet article compile contre Mathlib sans \lean{sorry}
et avec une garde d'axiomes qui fait échouer la construction en cas de dérive
(\S\ref{sec:verification}) ; rien ici n'est affirmé vrai parce qu'un modèle l'a dit.

\textbf{Ce que l'auteur a fait.} Xavier Callens a fixé l'orientation de la recherche, choisi quoi
formaliser et dans quel ordre, jugé ce que signifient les résultats et s'ils méritent d'être
publiés, et est responsable de chaque affirmation de ce document, y compris ses erreurs. Cette
responsabilité a été exercée concrètement, non nominalement, dans l'histoire même de cet article :
une version antérieure revendiquait une antériorité pour cette formalisation, revendication qui
s'est révélée fausse, a été rétractée et remplacée par l'attribution consignée au
\S\ref{sec:related} ; une tentative ultérieure d'étendre $W_N$ à une dualité physique
(\S\ref{sec:tduality}) s'est révélée, lors d'une relecture contradictoire, se réduire à un
renommage assorti d'une référence physique non soutenue, et est rapportée ici comme un résultat
négatif plutôt que close comme un théorème. Le test même de Tao est pertinent : un résultat que son
auteur ne peut ni expliquer ni défendre, vérifié formellement ou non, n'est pas prêt à être déclaré
complet~\cite{Tao2026AI}. Le contenu mathématique de $W_N$ --- ce qu'il est, pourquoi il normalise
$\Gamma_0(N)$, et pourquoi les deux revendications retirées ci-dessus étaient en fait erronées ---
est quelque chose que l'auteur peut expliquer sans assistance supplémentaire de l'IA ; les termes de
preuve Lean qui le certifient ont été rédigés pour l'essentiel par l'IA et ne sont dignes de
confiance que parce que le noyau, non l'auteur, les vérifie.

\begin{thebibliography}{99}
\bibitem{AtkinLehner1970} A.~O.~L.~Atkin, J.~Lehner, \textit{Hecke operators on $\Gamma_0(m)$}, Math.\ Ann.\ \textbf{185} (1970), 134--160. \texttt{doi:10.1007/BF01359701}.
\bibitem{Mathlib2020} The mathlib Community, \textit{The Lean mathematical library}, Proceedings of the 9th ACM SIGPLAN International Conference on Certified Programs and Proofs (CPP), 2020.
\bibitem{Buzzard2021} K.~Buzzard, \textit{Formalising mathematics}, Math.\ Proc.\ Cambridge Philos.\ Soc., 2021.
\bibitem{DiamondShurman2005} F.~Diamond, J.~Shurman, \textit{A First Course in Modular Forms}, Graduate Texts in Mathematics \textbf{228}, Springer (2005).
\bibitem{Miyake2006} T.~Miyake, \textit{Modular Forms}, Springer Monographs in Mathematics (2006).
\bibitem{Ono2004} K.~Ono, \textit{The Web of Modularity: Arithmetic of the Coefficients of Modular Forms and $q$-series}, CBMS Regional Conference Series in Mathematics \textbf{102}, AMS (2004). \texttt{doi:10.1090/cbms/102}.
\bibitem{Martin1996} Y.~Martin, \textit{Multiplicative eta-quotients}, Trans.\ Amer.\ Math.\ Soc.\ \textbf{348} (1996), 4825--4856.
\bibitem{GPR1994} A.~Giveon, M.~Porrati, E.~Rabinovici, \textit{Target Space Duality in String Theory}, Phys.\ Rept.\ \textbf{244} (1994), 77--202. \texttt{arXiv:hep-th/9401139}.
\bibitem{PerssonVolpato2015} D.~Persson, R.~Volpato, \textit{Fricke S-duality in CHL models}, JHEP \textbf{12} (2015) 156. \texttt{arXiv:1504.07260}.
\bibitem{FLT2026} Anthropic, \textit{A formal proof of Fermat's Last Theorem in Lean 4}, artefact logiciel, \url{https://github.com/anthropics/fermats-last-theorem}, 2026.
\bibitem{Prove2Me2026} S.~Chen, K.~Marwaha, X.~Lu, H.~Yuen, T.~Peng, \textit{Prove2Me: An open collaborative platform for scaling math formalization}, arXiv:2608.28433, 2026.
\bibitem{Tao2026AI} T.~Tao, \textit{Mathematics in the age of AI}, essai fondé sur une conférence
au Congrès international des mathématiciens, juillet 2026, soumis aux Actes du CIM 2026.
\texttt{arXiv:2608.16753}.
\bibitem{Leiden2026} \textit{La Déclaration de Leyde sur l'intelligence artificielle et les
mathématiques}, 2 juin 2026, endossée par l'Union mathématique internationale.
\url{https://leidendeclaration.ai/}, \texttt{doi:10.5281/zenodo.20302944}.
\bibitem{TaoEstimates2025} T.~Tao, \textit{estimates: a lightweight proof assistant for inequalities and asymptotic estimates}, logiciel, \url{https://github.com/teorth/estimates}, 2025.
\bibitem{deMoura2021} L.~de~Moura et al., \textit{The Lean 4 Theorem Prover and Programming Language}, CADE-28, LNCS \textbf{12699}, Springer (2021).
\end{thebibliography}

\end{document}
